% EEE3096S Practical 3: replace every blank with your own measured results.
% Place full12.png and coarse6.png alongside this file, or in docs/ when
% compiling from the practical root. Missing images produce safe placeholders.
% Compile from this directory: tectonic report_template.tex
% Keep this report to the measurements, two captures, and one page of explanation.
\documentclass[11pt,a4paper]{article}
\usepackage{lmodern}
\usepackage[margin=20mm]{geometry}
\usepackage{amsmath}
\usepackage{booktabs}
\usepackage{tabularx}
\usepackage{graphicx}
\usepackage[hidelinks]{hyperref}
\setlength{\parindent}{0pt}
\setlength{\parskip}{5pt}
\renewcommand{\arraystretch}{1.4}
\newcommand{\blank}[1][27mm]{\rule{#1}{0.4pt}}
\newcommand{\captureplaceholder}[2]{%
  \fbox{\parbox[c][0.26\textheight][c]{\dimexpr\linewidth-2\fboxsep-2\fboxrule\relax}{%
    \centering #2\par\medskip
    Replace with your own oscilloscope capture.\par
    Save it as \texttt{\detokenize{#1}}.}}}
\newcommand{\capture}[2]{%
  \IfFileExists{#1}{%
    \includegraphics[width=\linewidth,height=0.26\textheight,keepaspectratio]{#1}}{%
    \IfFileExists{docs/#1}{%
      \includegraphics[width=\linewidth,height=0.26\textheight,keepaspectratio]{docs/#1}}{%
      \captureplaceholder{#1}{#2}}}}

\begin{document}
\begin{center}
  {\Large\bfseries EEE3096S Practical 3 --- Voltage Mirror}\par
  Student report template
\end{center}
Names/student numbers: \blank[105mm]\par
Date: \blank[35mm]\qquad Build configuration: \blank[45mm]

\textbf{Assessment: live demonstration 80 marks; report 20 marks.}
Enter your own measurements. Suggested input levels are setup targets, not results.
Submit the four completed functions with this report.

\section*{1. Measurements and predictions (8 report marks)}
Measured analogue reference $V_R$: \blank\ V\par
DMM range/resolution: \blank[75mm]\par
Nominal ADC/DAC interval $V_R/4096$: \blank\ mV\par
Predicted coarse output step $V_R/64$: \blank\ mV

Measure PA5 and PA4 independently with the DMM. Record the actual LCD codes.
Calculate the nominal DAC output $D V_R/4096$ from each recorded command.

{\small
\begin{tabularx}{\linewidth}{@{}l l X r r X X@{}}
\toprule
Input target & Mode & Measured PA5 (V) & $C$ & $D$ & Predicted PA4 (V) & Measured PA4 (V) \\
\midrule
$\sim$0.800 V & FULL12  & & & & & \\
$\sim$0.800 V & COARSE6 & & & & & \\
$\sim$1.650 V & FULL12  & & & & & \\
$\sim$1.650 V & COARSE6 & & & & & \\
$\sim$2.500 V & FULL12  & & & & & \\
$\sim$2.500 V & COARSE6 & & & & & \\
\bottomrule
\end{tabularx}}

Generator: \blank[18mm]\ Hz; \blank[18mm]\ Vpp; \blank[18mm]\ V offset\par
Verified PA5 minimum/maximum: \blank\ / \blank\ V\par
Measured adjacent coarse step: \blank\ mV\par
Measured minus predicted step: \blank\ mV\par
Percentage difference $100(\Delta V_{\rm measured}-\Delta V_{\rm predicted})/
\Delta V_{\rm predicted}$: \blank[18mm]\ \%

\newpage
\section*{2. Two real scope captures (6 report marks)}
Label PA5 input and PA4 output, modes, voltage scales and time scales.
Use your own captures; simulated illustrations are not experimental evidence.

\begin{center}
\capture{full12.png}{FULL12: input and output over approximately one triangle cycle}
\end{center}
\textbf{Capture 1 --- FULL12.} Show both traces over approximately one cycle.
The small individual 12-bit steps need not be resolved.

\begin{center}
\capture{coarse6.png}{COARSE6: adjacent settled plateaus and voltage cursors}
\end{center}
\textbf{Capture 2 --- COARSE6.} Show both channels, adjacent settled output
plateaus and the cursor measurement used for your step comparison.

\newpage
\section*{3. Explanation (6 report marks; at most one page)}
Replace the prompts and blank spaces with concise answers supported by your
measurements. Keep the whole written explanation on this page.

\textbf{Mapping.} State your FULL12 and COARSE6 relationships. How many distinct
coarse commands exist, what is their spacing, and why is the maximum 4032?
\vspace{18mm}

\textbf{Measured steps.} Compare the measured coarse step with $V_R/64$.
Use your readings and instrument limits to discuss any difference.
\vspace{18mm}

\textbf{Resolution and accuracy.} Explain why a smooth FULL12 trace does not
establish continuous voltage resolution, and why 12-bit resolution does not
guarantee voltage accuracy. Both hardware converters stay at 12 bits.
\vspace{18mm}

\textbf{One DC result.} Compare $C V_R/4096$ with measured input, and
$D V_R/4096$ with measured output. Explain why copying the code does not
guarantee identical physical voltages.
\vspace{18mm}

\textit{Live demonstration: prepare all three DC levels; show one safe
tutor-selected input and both triangle modes. Explain the actual codes,
$V_R/64$ step prediction, ADC/DAC register fields, bounded wait and W1C flags.}
\end{document}
